\appendix

\section{Attribution of Assumptions (Which Proof Consumes What)}\label{app:A}

This appendix tabulates where the proofs in the text consume inputs
1--8 of Theorem~\ref{thm:central}. The content of the claims is
governed by the corresponding sections of the text; this appendix is an
index of attribution.

\subsection{Comparison core (\S\ref{sec:53}--\S\ref{sec:55})}\label{app:A1}

\begin{itemize}
\item Consumed: the cover intersection diagram; the two coefficient
presheaves (assumptions 1--2); the generator map and the two
completeness conditions (assumptions 3--4); the local-state data
(assumptions 5--7, used in the derivation of soundness).
\item Not used: the global sheaf condition, the finite enumeration of
the cover, displayed equation fulfillment.
\end{itemize}

\subsection{Actual gluing (\S\ref{sec:56})}\label{app:A2}

\begin{itemize}
\item What Lemma~5.2A consumes is the monomorphism property of the
cover and, of assumption~8, the subsingleton clause for $\Psem(V)$
(the proof consumes only the pairwise-overlap part). The lemma is used
in the completion of the corrected family in the forward proof and in
the gluing step of the equation-side globalization.
\item In addition, actual gluing uses the true sheaf condition of
(iii) (the sheaf condition and membership in the topology). The
$\Psem(V)$ subsingleton clause is supplied as the same proposition as
condition~(4) of the true sheaf condition.
\end{itemize}

\subsection{Unconsumed clauses}\label{app:A3}

Of assumption~8, the subsingleton clause for $\PE(V)$ and the
coefficient-vanishing clause $\Msem(V)=\QE(V)=0$ are consumed by no
proof in this paper. The complexes run over nonempty intersections
only (\S\ref{sec:41}), and what the equation-side globalization (end
of \S\ref{sec:56}) uses for empty overlaps is also the
$\Psem$-side subsingleton clause (Lemma~5.2A). These unconsumed
clauses are retained in order to fix the input side (assumption~8) in
a form that includes the full identification with the ordered
\v{C}ech complex of \S\ref{sec:41} (outside the scope of this paper).

\subsection{Consumption in the Lean formalization (thin sites)}\label{app:A4}

The Lean formalization is carried out over a site whose context
category is thin (\S\ref{sec:57}). By thinness --- the agreement of
parallel morphisms --- every morphism in this model is automatically a
monomorphism, and the treatment of self-overlaps and reversed overlaps
in Lemma~5.2A is also settled automatically by the agreement of
parallel morphisms. Consequently the formalized counterpart of
Lemma~5.2A, \code{SiteStateData.matchingFamily\_iff}, consumes the
selected pairwise overlaps and their lift data, but does not consume
monomorphy as an explicit hypothesis (the commutativity and uniqueness
of pullbacks follow from thinness). The gluing argument over general
$\ArchCtx(X)$ is \code{unported} (Section~\ref{sec:6}).

\section{Lean Source Correspondence (Declarations, Labels, Source Files)}\label{app:B}

This appendix fixes the organization, within the release tag, of the
Lean sources that are the primary evidence for the status of
Section~\ref{sec:6}. The SAGA theorem chain is located under
\code{Formal/AG/SemanticRepair/Saga/}. The kernel axiom audit is
implemented by \code{\#assert\_standard\_axioms\_only} in
\code{Formal/AG/AxiomAudit.lean} (an allowlist check of all registered
declarations), which the release CI executes as
\code{lake env lean Formal/AG/AxiomAudit.lean}.

The docstrings of the Lean sources use labels from a numbering series
distinct from this paper, originating in internal organization during
development (\code{X.Theorem 1.1} and so on). The theorem
correspondence with this paper is given completely by the status table
of Section~\ref{sec:6}; the labels serve for cross-checking when
browsing the Lean sources. The correspondence with the declarations of
Section~\ref{sec:6} is as follows (declarations without labels are
marked ``---''; source files are relative to \code{Saga/}).

\begin{table}[htbp]
\centering
\scriptsize
\begin{tabular}{@{}p{0.46\textwidth}p{0.22\textwidth}p{0.24\textwidth}@{}}
\toprule
Lean declaration & Lean-side label & Source file \\
\midrule
\code{SagaEquationPacket.sagaCentralTheorem} & X.Theorem 1.1 &
\code{TrueSheafDescent.lean} \\
\addlinespace
\code{PrimaryStateCorrespondence.}\allowbreak\code{relationSound\_of\_stateCorrespondence} &
--- & \code{Exactness.lean} \\
\addlinespace
\code{SagaEquationPacket.phiEquiv} (generic version:
\code{PrimaryCoefficientCorrespondence.phiEquiv}, X.Theorem 6.3 /
Corollary 6.7) & as listed & \code{EquationRealization.lean},
\code{Exactness.lean} \\
\addlinespace
\code{ExactnessFixtures.soundness\_failure} /
\code{completeness\_failure} / \code{generation\_failure} &
X.Example 6.6 & \code{Exactness.lean} \\
\addlinespace
\code{kappa1\_delta0}, \code{kappa2\_delta1} & X.Theorem 7.2 &
\code{KappaComparison.lean} \\
\addlinespace
\code{SagaEquationPacket.kappaStarAddEquiv} (generic version:
\code{kappaH1AddEquiv}, X.Theorem 7.4) & as listed &
\code{KappaComparison.lean} \\
\addlinespace
\code{SagaEquationPacket.residual\_correspondence\_class},
\code{betaAligned\_residual} & X.Theorem 7.5 &
\code{KappaComparison.lean} \\
\addlinespace
\code{SagaEquationPacket.globalRepair\_nonempty\_iff},
\code{sagaGroundedGluing} & X.Theorem 8.2 &
\code{TrueSheafDescent.lean} \\
\addlinespace
\code{SiteStateData.matchingFamily\_iff} & --- &
\code{OrderedComparison.lean} \\
\addlinespace
\code{CircleWitness.semanticResidualClass\_ne\_zero},
\code{circle\_nonzero\_class\_transfer} & X.Example 10.2 / Appendix
B.9 & \code{CircleWitness.lean} \\
\addlinespace
\code{DescentWitness.descentTrueSheaf},
\code{descent\_sagaGroundedGluing} & --- &
\code{DescentWitness.lean} \\
\bottomrule
\end{tabular}
\caption{Correspondence between declarations, Lean-side labels, and
source files.}
\label{tab:leandecl}
\end{table}

The correspondence between the building blocks of \S\ref{sec:61} and
the source files:

\begin{table}[htbp]
\centering
\scriptsize
\begin{tabular}{@{}p{0.55\textwidth}p{0.36\textwidth}@{}}
\toprule
Building block (\S\ref{sec:61}) & Source file \\
\midrule
cover, intersection diagram, three-term \v{C}ech complex &
\code{Cover.lean}, \code{CechThreeTerm.lean} \\
semantic repair presentation and $\Msem$ & \code{Presentation.lean} \\
affine repair system, selected atlases, residuals &
\code{RepairTorsor.lean}, \code{EquationLift.lean} \\
soundness derivation, exactness, coefficient isomorphism, finite
counterexamples & \code{Exactness.lean} \\
equation-side realization / production &
\code{EquationRealization.lean}, \code{EquationProduction.lean} \\
cochain comparison $\kappa$, $H^1$ isomorphism, residual
correspondence & \code{KappaComparison.lean} \\
true sheaf descent, final bundling of the central theorem &
\code{TrueSheafDescent.lean} \\
finite witnesses of the zero and nonzero cases &
\code{CircleWitness.lean}, \code{DescentWitness.lean} \\
poset site realization, ordered \v{C}ech model comparison bridge &
\code{OrderedComparison.lean}, \code{PartIVBridge.lean} \\
\bottomrule
\end{tabular}
\caption{Correspondence between the building blocks of
\S\ref{sec:61} and source files.}
\label{tab:leanfiles}
\end{table}
